\section{Observations, affine Gaussian shapes, and isolation}
\label{sec:preparation}

We prepare the critical Bessel height law for the local height maps.
The preparation has two parts.  First, an entire cutoff changes the
finite-volume quantities under consideration by an exponentially small
amount.  Second, an exact isolation expansion expresses the cutoff
interaction in terms of tile errors, restoration errors, and exponentially
small Gaussian links.  The affine Gaussian estimates below allow the same
construction in the sector with one circulation.  The next section proves
the analytic norm estimates needed to enter the maps.

\subsection{Observation cells and the cutoff}
\label{prep:geometry}

Set $\alpha=\sqrt{8\pi}$, and use the nearest-neighbor polynomial
covariances and subset algebra specified in the inputs.  Thus $A$ is the
unit-edge height Laplacian, and the covariance integrated below side $m$
is
\begin{equation}\label{prep:entry-covariance}
 C_{<m}=A^{-1}(1-\mathsf P_m),
\end{equation}
with polynomial continuation on constants and with the spectral-argument
convention fixed there.  There is no additional independent-variance
factor.  We use the same polynomials on free dual squares.  The remaining
reference field has a common constant uniform modulo $2\pi/\alpha$.
Statements about the plane concern local formulas; no probability law on
absolute plane heights is required.  The block collars, analytic tests,
regulators, and compatibility rules are those of \cite[Sections
9, 8, and 4]{B}.
The tile parameter $R$ below is distinct from the fixed block connection
radius $\rho$.

For a dyadic integer $s$, partition the microscopic graph into square
\emph{observation cells} of side $s$.  Write $B_s$ for cell averaging, and
introduce independent standard Gaussian variables $e_c$, one per cell.
The observations are
\begin{equation}\label{prep:observations}
                  y=B_s h+e.
\end{equation}
All cuts and clippings in this section retain whole observation cells.
In particular a retained observation face consists of all the microscopic
bonds joining the two cells on its sides.

In observation-cell units use the tile geometry of
\cite[Theorem 9.1 and Subsection 9.2]{B}.
Start with squares of side $R$, remove the square of side $R/2$ centered
at each grid vertex from the incident squares, and make each removed
square a vertex tile.  The remaining parts of the original squares and
the vertex tiles partition the cells.  Here $R$ is a large dyadic integer,
$P=R^5$, $q$ is a fixed dyadic period parameter large enough for all fixed
collars, and
\begin{equation}\label{prep:scales}
                         m=qRs.
\end{equation}
The block factor $L$ is a sufficiently large power of two.  We use square
and torus sides compatible with the grids, at least $sP$, and divisible by
$sP$.  Free-wall tiles are the allowed box clips of this construction;
their cuts lie between cell layers.

Choose the common square-symmetry origins at grid vertices, which are
elementary-square centers in site coordinates.  In a pierced square put
the source at the central grid vertex and take the integral connection
$l$ along a grid ray.  Thus $l$ vanishes inside each observation cell and
has a constant value $l_f$ on the microscopic edges crossing an oriented
observation face $f=(c,c')$.  The source lies in the central vertex tile,
even though four observation cells meet there.  Tori in this proof are
unpierced and carry ordinary single-valued heights.

On every observation face retained in the graph insert
\begin{equation}\label{prep:cutoff-definition}
 \chi_T(d_f)=\int_{-T}^{T}\frac{1}{\sqrt{2\pi}}
                   e^{-(d_f-x)^2/2}\,dx,
 \qquad d_f=y_{c'}-y_c+l_f,\qquad T=\sqrt{s}.
\end{equation}
The function $\chi_T$ is entire.  For real arguments it lies in $(0,1)$,
and $1-\chi_T(d)\le C e^{-cT^2}$ for $|d|\le T/2$.  Its role is to give
the observation law a quadratic large-field reserve on the finite scale
intervals used later.

\begin{lemma}[Finite-volume cutoff error]\label{prep:cutoff}
On a compatible torus or free square of side $M$, the absolute change in
the partition sum caused by the product of
\eqref{prep:cutoff-definition}, divided by the centered partition
sum without cutoffs,
is at most
\begin{equation}\label{eq:cutoff-error}
                         C(M/s)^2 e^{-c_1s}.
\end{equation}
The assertion includes a free square with the integral connection $l$.
On a torus it also holds, with a smaller $c_1>0$, after inserting a
bounded test, a fixed power of a smooth mean-free smear, or a fixed real
or complex exponential test of such smears.
\end{lemma}

\begin{proof}
Consider adjacent cells $c,c'$.  Send the uniform mass on $c$ to its
translate in $c'$ along parallel normal lattice paths.  Every path has
length $s$.  Giving each starting site mass $s^{-2}$ produces an edge
flow $a_f$ with
\begin{equation}\label{prep:two-cell-flow}
 \langle a_f,\partial h+l\rangle
       =(B_sh)_{c'}-(B_sh)_c+l_f,
 \qquad \max_e|a_f(e)|\le C/s,
 \qquad \sum_e a_f(e)^2\le C.
\end{equation}
Indeed at most $s$ starting sites use any given edge, and there are
$O(s^2)$ edges in the two cells.  Each path crosses the common face
once, so its affine offset supplies precisely $l_f$.  Consequently
$d_f=\langle a_f,\partial h+l\rangle+e_{c'}-e_c$.

Apply the tilted shifted-to-centered comparison of
\cite[Lemma 2.2]{X}, first on the independent
Gaussian chains in series.  Release their endpoint, cycle, and period
constraints, and release the independent observation noises.  After the
fixed-graph series limit the log transform of one primary increment at
argument $u$ is $b_c(\cosh(2\pi u)-1)$.  If $|w|\le c\sqrt{s}$, all
arguments $wa_f(e)$ are in a fixed bounded interval.  Therefore
\[
 \sum_e b_c\bigl(\cosh(2\pi wa_f(e))-1\bigr)
       +\tfrac12w^2\norm{\mathbf 1_{c'}-\mathbf 1_c}^2
       \le Cw^2.
\]
The comparison gives this bound for the affine numerator normalized by
the centered partition sum as well.  Choosing $w$ proportional to
$\sqrt{s}$ and applying the exponential Markov inequality to both signs
shows that the centered-normalized mass of $|d_f|>T/2$ is at most
$Ce^{-cs}$.

On $|d_f|\le T/2$ use the Gaussian tail bound for $1-\chi_T(d_f)$;
on the complement use $0\le1-\chi_T\le1$.  Since
$1-\prod_f\chi_T(d_f)\le\sum_f(1-\chi_T(d_f))$ and there are
$O((M/s)^2)$ faces, this proves \eqref{eq:cutoff-error}.
For the asserted torus insertions, H\"older's inequality separates the
rare-event factor from the test.  The latter has the required higher
moments and exponential moments by
\cite[Lemma 2.3]{X}.  A fixed complex
exponential is handled by its real part in the absolute-value bound.
The H\"older exponent only decreases the positive rate in
\eqref{eq:cutoff-error}.
\end{proof}

For the pierced square, the input lower bound
$a_n\ge c n^{-1/2}$ converts \eqref{eq:cutoff-error} into a relative
error bound.  Its use here requires no information about the eventual
normalizing amplitude; the bound comes from
\cite[Equations (18)--(19)]{S} and the fixed-spin duality of its
Section~2.  In particular a polynomial loss when dividing by the magnetic
partition ratio is harmless when $s$ is chosen on the finite histories
below.

\subsection{The affine observation identity and local lifts}
\label{prep:gaussian-identity}

Let $G$ be a graph obtained by the cell cuts just described.  Write
$\mathcal F(G)$ for its retained observation faces.  Free constants are
allowed on every connected component.  Define the Gaussian shape energy
by
\begin{equation}\label{eq:shape-energy}
 E_G^l(y)=\inf_v\left\{\alpha^{-2}|\partial v+l|_G^2
                                      +|B_sv-y|^2\right\}.
\end{equation}
Write $E_G=E_G^0$.  Let $Z_G^l(y)$ be the Bessel height sum with observation
penalty $\exp(-|B_sh-y|^2/2)$ and with all component constants summed.
No cutoff occurs in this definition.  More explicitly, apart from the
common normalization of the observation noise, it is
\[
 Z_G^l(y)=\sum_{h\in(2\pi\Z)^{V(G)}}
       \prod_{e\in E(G)}p_{b_c}\!\left(\frac{\partial h(e)+l(e)}{2\pi}\right)
                         e^{-|B_sh-y|^2/2}.
\]
Here $E(G)$ denotes the set of microscopic edges, each counted once;
the observation penalty makes the sum over each component constant
finite.  When subsequently integrating observations, one divides out
common translations, or uses one period for each free constant.

On the full connected graph, let $r$ minimize the unobserved affine
energy, and put $\eta_l=\partial r+l$.  The minimizer is determined up to
a constant and satisfies
\[
                         \partial^*\eta_l=0.
\]
For $v=r+u$, orthogonality gives
$|\partial v+l|^2=|\eta_l|^2+|\partial u|^2$.  Gaussian integration and
completion of the observation square consequently give the affine
Gaussian observation law
\begin{equation}\label{eq:gaussian-observation}
 y=B_sr+B_s\alpha\psi+e,
 \qquad \operatorname{cov}'(\psi)=A^+.
\end{equation}
The prime omits constants.  The constant of $\psi$ is uniform modulo
$2\pi/\alpha$, and $e$ has the independent standard Gaussian law already
specified.  This calculation may equivalently be made by first summing
height translations and then integrating $y$, with the common translation
divided out.

For clarity, define the interaction on every cut graph by
\begin{equation}\label{eq:interaction}
 J_G^l(y)=e^{E_G^l(y)/2}Z_G^l(y)
                       \prod_{f\in\mathcal F(G)}\chi_T(d_f).
\end{equation}
On the full graph suppress $G$.  Relative to the unpierced Gaussian
integral, evaluating $J^l$ under \eqref{eq:gaussian-observation} introduces
the scalar
\begin{equation}\label{prep:magnetic-scalar}
                         e^{-|\eta_l|^2/(2\alpha^2)}.
\end{equation}
To see that this is the complete scalar, the Gaussian integral with
fixed observations equals a common determinant factor times
$e^{-E^l(y)/2}$.  Translating by $r$ makes its determinant exactly the
unpierced determinant and leaves the constant energy
$|\eta_l|^2/\alpha^2$.  Multiplication by \eqref{eq:interaction} cancels
the observation shape.  Thus no additional factor involving a chosen
local harmonic lift is present.

We record the lift estimates that determine how the eventual activities
copy between the plane and a square.  They are consequences of
\cite[Lemmas 4.2 and 5.2, and
the duality in Lemma 5.1]{B}, with the source period rescaled
to $2\pi$ in height units.

Extend the chosen integral ray connection $l$ to the plane.  Let
$\eta_\infty$ be the divergence-free infinite-plane connection obtained
as the local limit of the square minimizing connections, equivalently
the rotated infinite-plane Green gradient normalized to circulation
$2\pi$ at the source and zero circulation elsewhere.  Its energy grows
logarithmically with the radius; no finite total plane energy is assumed.
Since $\eta_\infty-l$ has zero circulation around every lattice plaquette,
path integration defines a single-valued real height $r_\infty$ by
\[
                         \partial r_\infty=\eta_\infty-l.
\]
Use the same ray convention for $r$ and $r_\infty$, and align their
additive constants at one fixed height vertex adjacent to the mark,
chosen consistently with the grids.

\begin{lemma}[Local lifts and comparison of squares]\label{prep:lifts}
Let $r$ be the square minimizer at side $M$ and $r_\infty$ the plane
height just defined.
If an axis rectangle and its fixed collar avoid the source, choose an
integral gauge potential $g$ there with $\partial g=l$.  The local real
lift is $\lambda=(r+g)/\alpha$, and
$\partial\lambda=\eta_l/\alpha$.  At interior supports meeting the source,
the translation from the pierced-plane calculation to the square is
\begin{equation}\label{eq:lift-difference}
                         \gamma_M=(r-r_\infty)/\alpha.
\end{equation}
In lattice units, for the fixed derivative orders used in the maps,
\[
 |\nabla^a\eta_l(x)|\le C_a(1+\operatorname{dist}(x,\text{source}))^{-1-a}.
\]
In an inner fixed fraction of the square,
\begin{equation}\label{eq:lift-derivatives}
                 |\nabla^a\gamma_M|\le C_aM^{-a},
                 \qquad 1\le a\le5.
\end{equation}
The connection estimate includes the reflected stencils at free walls.
The affine contribution $(B_sr)_{c'}-(B_sr)_c+l_f$ to every retained
observation difference is bounded uniformly in $s$ and $M$.
\end{lemma}

\begin{proof}
The minimizing connection is the rotated killed-Green gradient under the
crossing correspondence.  Away from its source its circulation vanishes,
so path integration gives the lift on the stated rectangle.  The
positive-order Green estimates in the cited lemmas give the bound on
$\eta_l$.  At a free wall, odd reflection of the killed potential makes
the rotated normal component odd and the tangential component even,
which supplies the required reflected lift.

In a fixed interior fraction of the box, all finite-range polynomial
pieces of side below $cM$ coincide exactly with their plane pieces.
For a derivative of order $a\ge1$, each remaining piece is bounded by
$C_at^{-a}$ at its side $t$.  Summation over dyadic $t\ge cM$ gives
\eqref{eq:lift-derivatives}; the smoothed terminal piece has the same
bound.  Only positive-order differences of $r-r_\infty$ are used, so
the additive normalization does not affect this estimate.

Finally use the two-cell flow from \eqref{prep:two-cell-flow} with
$\eta_l$ in place of $\partial h+l$.  For cells within distance $O(s)$ of
the mark, its absolute value is bounded by
\[
 \frac Cs\sum_{e:\operatorname{dist}(e,\text{source})\le Cs}
              \frac{1}{1+\operatorname{dist}(e,\text{source})}\le C.
\]
For a two-cell pair farther away, the same bound is
$Cs/(1+d)\le C$, where $d\ge cs$ is its distance from the mark.
This includes the few faces adjacent to the pierced vertex.
\end{proof}

For later copying, a support is called \emph{shift-good} when its spanning
axis rectangle, enlarged by $Hm$ for a sufficiently large fixed $H$,
misses the mark.  On such a support the connection gauges away and the
ordinary calculation is evaluated at $\psi+\lambda$.  At an interior
support near the mark, the pierced-plane calculation is evaluated at
$\psi+\gamma_M$.  The spanning rectangle in this definition will include
the numerical dependencies of links, as specified below.  Thus these
statements apply to the entire calculation on a support, including an
insertion, and not just to the fields explicitly occurring in a factor.

\subsection{A real large-observation bound}
\label{prep:large-observations}

The cutoff yields a uniform strict saving in face differences.  We first
prove the estimate for the height sum without the cutoff.  It is useful
on every open patch, including a clipped patch carrying the connection.

\begin{lemma}[Large observations]\label{prep:large-field}
For every preparation graph $G$ and every real observation vector $y$,
\begin{equation}\label{eq:large-observation}
 \frac{Z_G^l(y)}{Z_G^0(0)}
       \le\exp\left\{-c\sum_{\substack{f\in\mathcal F(G)\\|d_f|\le2T}}
                                                    d_f^2\right\}.
\end{equation}
Consequently, with a possibly smaller constant $c>0$ independent of
$R,P,s$,
\begin{equation}\label{prep:face-saving}
 \frac{Z_G^l(y)}{Z_G^0(0)}
             \prod_{f\in\mathcal F(G)}\chi_T(d_f)
                  \le e^{-c\sum_{f\in\mathcal F(G)}d_f^2}.
\end{equation}
\end{lemma}

\begin{proof}
For each face with $|d_f|\le2T$, insert a real tilt
$w_f=c_0d_f$ on the two-cell flow and on the corresponding noise
difference.  At the prescribed observation $y$, the total statistic is
the deterministic number $\sum_fw_fd_f$.  The flows have bounded overlap:
each is supported on its two cells and each cell meets at most four
faces.  Hence their combined microscopic coefficients $u_e$ and noise
coefficients $v_c$ satisfy
\[
 \max_e|u_e|\le C\max_f|w_f|/s,
 \qquad \sum_eu_e^2+\sum_cv_c^2\le C\sum_fw_f^2.
\]

We spell out the affine comparison used here.  On the Gaussian-chain
lattice, impose endpoint matches and observations as squared equality
penalties.  Put each prescribed offset, including $l$ and the observation
value, into an additional coordinate fixed to that value in the
numerator and fixed to zero in the denominator, with no original energy
on that coordinate.  The constraints are then the same positive
semidefinite additions in the affine and centered coordinates.  Tilt the
affine edge coordinates and the noise variables themselves.
By \cite[Lemma 2.2(ii)]{X}, adding these common
precisions decreases the tilted-to-centered ratio.  Continuous prescribed
observations follow first from finite quadratic penalties and then from
the density limit.  Independent continuous Gaussian variables are
allowed in the same comparison, or follow by lattice approximation.

To pass between absolute site heights and relative coordinates, one may
temporarily retain independent site heights with auxiliary positive
precision and send that precision to zero.  Released tilts involve only
edge and noise variables, so the auxiliary factors cancel.  The
constrained quantities in the limit are exactly $Z_G^l(y)$ and $Z_G^0(0)$,
with constants summed as defined above.  All these limits are at a fixed
primary graph.

Upon releasing the constraints and taking the series limit, the log
transform is at most $C\sum_fw_f^2$.  Indeed
$\max|w_f|\le2c_0\sqrt{s}$ keeps the edge arguments uniformly bounded,
and the noise transform is quadratic without restriction.  Therefore
\[
 e^{\sum_fw_fd_f}\frac{Z_G^l(y)}{Z_G^0(0)}
                       \le e^{C\sum_fw_f^2}.
\]
Choose the fixed $c_0>0$ so small that $c_0-Cc_0^2>0$; this proves
\eqref{eq:large-observation}.  For $|d|>2T$, the Gaussian tail in
\eqref{prep:cutoff-definition} gives $\chi_T(d)\le e^{-c d^2}$.
Combining the two ranges proves \eqref{prep:face-saving}.
\end{proof}

\subsection{Localization of the affine shape energy}
\label{prep:gaussian-localization}

The Gaussian part of the isolation expansion compares energies on
different cut graphs.  We need these comparisons to be localized in
observation-cell distance with constants independent of $s$.  A nonzero
circulation creates a logarithmic energy near a cell junction.  The next
lemma isolates that energy in an exactly local term; the remaining
matrix has uniform exponential bounds.

For an observation junction $z$ at which all four faces are retained,
let $\kappa_z(d)$ be their oriented circulation.  No circulation term is
assigned to a junction whose retained sectors have tree adjacency.
For data $d_f=y_{c'}-y_c+l_f$, all $\kappa_z(d)$ vanish except possibly at
the source, where its absolute value is $2\pi$ if the full star is
retained.

\begin{lemma}[Affine shape localization]\label{prep:localization}
On each open preparation graph there are a quadratic form $L_G$ with a
fixed-radius stencil in the faces, and a symmetric matrix $N_G$ on those
faces, such that
\begin{equation}\label{prep:shape-decomposition}
 E_G^l(y)=L_G(d)+\sum_{f,g\in\mathcal F(G)}N_G(f,g)d_fd_g,
\end{equation}
where
\begin{equation}\label{prep:trial-energy}
 0\le L_G(d)\le C\sum_fd_f^2
                       +C\log(2s)\sum_z\kappa_z(d)^2.
\end{equation}
The forms and matrices may be defined on arbitrary face arrays $d$,
without imposing a curl constraint.  Their nonlocal entries satisfy
\begin{equation}\label{prep:matrix-decay}
                         |N_G(f,g)|\le Ce^{-c\operatorname{dist}(f,g)}.
\end{equation}
Their dependence on the graph is exponentially localized as well.  In
particular, if faces in a set $D$ are deleted, extend both matrices by
zero on missing faces.  With a fixed stencil enlargement understood in
the distances,
\begin{equation}\label{prep:matrix-cut}
 |(N_G-N_{G\setminus D})(f,g)|
 \le C\exp\left\{-c\left(\operatorname{dist}(f,g)
       +\min\{\operatorname{dist}(f,D),\operatorname{dist}(g,D)\}\right)
                                                   \right\}.
\end{equation}
The local stencils agree wherever the graphs agree.  For the actual
one-source data this implies
\begin{equation}\label{eq:affine-shape-bound}
 E_G^l(y)\le C\sum_{f\in\mathcal F(G)}d_f^2
                  +C\log(2s)\mathbf 1_{\{\text{pierced star present}\}}.
\end{equation}
\end{lemma}

\begin{proof}
\emph{Trial extensions.}
Subtract $y_c$ from the sites of each cell $c$.  If $u$ is the resulting
site field, then \eqref{eq:shape-energy} becomes
\[
 \inf_u\left\{\alpha^{-2}|\partial u+j(d)|_G^2+|B_su|^2\right\},
\]
where $j(d)$ equals $d_f$ on the microscopic edges crossing face $f$ and
is zero inside cells.  Construct a linear trial extension $t(d)$ using
only a fixed number of cells near each datum.  Away from face endpoints,
put opposite half-jumps on its two sides, with the signs canceling
$j(d)$, and damp these values into the cells at scale $s$.  The new
differences are bounded by $C/s$ times the data size.

At a vertex with tree adjacency, choose constants on its sectors to
cancel the retained jumps.  At a cyclic vertex, first do this for the
part with zero total circulation.  For the remaining circulation use,
inside a disk of radius $cs$, the exact lattice angular field: prescribe
the affine gradient to be $\kappa_z(d)/(2\pi)$ times the divergence-free
infinite-plane connection normalized to circulation $2\pi$, obtained by
rescaling \cite[Lemma 5.2]{B}.
Path integration constructs the sector fields because the prescribed
and original circulations agree.  For unit data its gradient is bounded
by $C/(1+d)$ at distance $d$ from the junction.

The constants can be chosen so that the values used for gluing at radii
comparable to $s$ are bounded.  Integrating along arcs in an annulus of
fixed radii ratio gives this assertion; successive inner dyadic annuli
give, if needed, the sufficient value bound
$C(1+\log^+(s/(1+d)))$.  In particular its cell average and its
$s^{-2}$-weighted squared site norm are bounded uniformly in $s$.
Glue the vertex fields, face strips, and zero cell-interior extensions
with a partition of unity at scale $s$.  For example, vertex cutoffs may
equal one within sup distance $s/10$ and vanish beyond $s/5$; the face
strips outside the smaller neighborhoods have width less than $s/10$
and central plateaus.  On the two sides of a joined face use matching
cutoffs, with differences $O(s^{-1})$.  At free cuts join only sectors
connected by retained faces.  Averaging the fixed stencils over square
symmetries preserves these properties.

Define
\[
 L_G(d)=\alpha^{-2}|\partial t(d)+j(d)|_G^2+|B_st(d)|^2.
\]
It is exactly local.  Its face-strip and gluing energies are bounded by
$C\sum_fd_f^2$.  In the inner cyclic neighborhoods, summing
$(1+d)^{-2}$ over lattice annuli costs $C\log(2s)$ times the squared
circulation.  This proves \eqref{prep:trial-energy}.

\emph{The correction source.}
Let $Q_G$ be the homogeneous massive observation form
\begin{equation}\label{prep:massive-form}
 Q_G(u)=\alpha^{-2}|\partial u|_G^2+|B_su|^2,
 \qquad Q_G(u,v)=\alpha^{-2}\langle\partial u,\partial v\rangle_G
                                      +\langle B_su,B_sv\rangle.
\end{equation}
Cell Poincar\'e gives, on any union of whole cells,
\begin{equation}\label{prep:cell-poincare}
 s^{-2}\sum_c\norm{u}_{\ell^2(c)}^2
       \le C\left(\sum_c\norm{\partial u}_{\ell^2(c)}^2
                                          +\sum_c|(B_su)_c|^2\right).
\end{equation}
This form controls component constants, including on cut graphs.

The linear functional produced by the trial field is
\[
 q_d(u)=\alpha^{-2}\langle\partial t(d)+j(d),\partial u\rangle_G
                                        +\langle B_st(d),B_su\rangle.
\]
Write $q_d=\sum_fd_fq_f$.  For each unit face datum, $q_f$ is supported
on a fixed number of nearby cells, and we claim
\begin{equation}\label{prep:bounded-source}
                         |q_f(u)|\le C Q_G(u;U_f)^{1/2}.
\end{equation}
Here $U_f$ is a fixed enlargement of those cells and $Q_G(u;U_f)$
includes their cells and incident retained faces.  The only issue is the
angular gradient, whose full squared norm has the logarithmic bound
just obtained.  Choose a cutoff $\xi$ equal to one on its inner disk and
supported in a slightly larger disk still inside the exact angular
region.  The connection $\eta_\infty$ there is divergence-free, so
\[
                     \langle\eta_\infty,\partial(\xi u)\rangle=0.
\]
Subtract this identity from its pairing with $\partial u$.  All surviving
terms lie in an annulus with radii comparable to $s$.  The connection has
bounded squared norm on that annulus.  Terms containing $\partial\xi$
are bounded using $|\partial\xi|\le C/s$ and
\eqref{prep:cell-poincare}.  The other trial pieces have bounded energy,
and $B_st$ is bounded as proved above.  Cauchy--Schwarz now gives
\eqref{prep:bounded-source}.  Thus the logarithmic core energy has not
entered the norm of the correction source.

Let $H_G$ be the positive operator associated with $Q_G$, and let $u_f$
solve $Q_G(u_f,v)=q_f(v)$.  Completing the square around $t(d)$ gives
\[
 E_G^l(y)=L_G(d)-q_d(H_G^{-1}q_d),
 \qquad N_G(f,g)=-q_f(H_G^{-1}q_g).
\]
This proves \eqref{prep:shape-decomposition}; it also shows that the
nonlocal matrix is the negative Gram matrix of bounded local sources.
In particular \eqref{eq:affine-shape-bound} already follows by using the
trial field and the stated circulation pattern.

\emph{Energy decay away from a source.}
We give the uniform localization argument because it is also needed for
graph comparison.  A solution $u_f$ is $Q_G$-harmonic outside $U_f$.
Choose sup-distance boxes in cell units and a cutoff $\zeta$ which is
zero inside one box, one beyond a larger box, and changes on a layer of
a fixed large number of cells.  Interpolate it on sites with
$|\partial\zeta|\le C/s$ and test the equation with $\zeta^2u_f$.
The form on cells and joining faces wholly outside the transition layer
is the corresponding outside energy.  Gradient cross terms and mass
terms in the transition region are bounded by its energy, with one
extra layer on either side, using Cauchy--Schwarz and
\eqref{prep:cell-poincare}.  Consequently, if $F(k)$ is the energy outside
a box at cell distance $k$, there is a fixed width $b_0$ for which
\[
 F(k+2b_0)\le C\bigl(F(k-b_0)-F(k+2b_0)\bigr).
\]
Rearranging gives a fixed fractional decrease every $3b_0$ layers.
The total energy is bounded by \eqref{prep:bounded-source}, so
\begin{equation}\label{prep:solution-decay}
                    Q_G(u_f;\operatorname{dist}(\cdot,U_f)\ge k)
                                          \le Ce^{-ck}.
\end{equation}
The proof is unchanged for arbitrary preparation cuts.  Intrinsic face
distance can also be used: assign separate cutoff values to disconnected
local sectors and interpolate along retained faces.  At a joined
junction the discrepancies are at most four cell steps, so the same
gradient bound applies.

Applying the local bound \eqref{prep:bounded-source} for $q_g$ to
$u_f$ and then \eqref{prep:solution-decay} proves
\eqref{prep:matrix-decay}, after decreasing $c$.

\emph{Comparison of two graphs.}
Suppose two graphs agree in a box of cell radius $k$ about $U_f$.
Multiply the solution in the first graph by a reversed cutoff, equal to
one near $U_f$ and zero beyond a transition layer halfway out.  Transplant
it to the second graph and put it equal to zero elsewhere.  Its equation
error against a unit-$Q_G$ test is bounded by $Ce^{-ck}$: pair the uncut
equation with the analogously cut test, and bound the resulting layer
terms by \eqref{prep:cell-poincare} and the square root of the layer
energy in \eqref{prep:solution-decay}.  Coercivity of $Q_G$ bounds the
difference from the second solution by the same quantity in energy
norm.  This proves exponentially local dependence of the solution and
of the matrix entries.

For a deletion set $D$, each of the two matrices has exponential decay
in the entry separation.  If both entries are far from $D$, their local
stencils agree.  When their separation is small compared with their
distance from $D$, compare the solutions in a common box as just done.
When the separation is larger, the separate matrix bounds already give
exponential decay in that distance.  Thus the difference has both a
separation bound and a bound in the minimum distance to $D$.  Taking
geometric means of these bounds proves the simultaneous estimate
\eqref{prep:matrix-cut}, with smaller exponential constants.  Missing
rows and columns cause no problem under the zero-extension convention.
\end{proof}

The same proof supplies the extra estimate used in isolation.  Compare
a true deletion with its calculation on an open patch having clearance
of order $R$ from artificial boundaries.  For entries near the deletion,
the graph-comparison argument has a common box of radius $cR$.  For
entries far from it, each deletion change is already exponentially
small by \eqref{prep:matrix-cut}.  Combining these bounds with separation
decay gives, for the error $\Delta_i$ at event $i$ with deletion region
$D_i$,
\begin{equation}\label{prep:event-mismatch}
 |\Delta_i(f,g)|\le
 C\exp\left\{-c\left(R+\operatorname{dist}(f,g)
       +\min\{\operatorname{dist}(f,D_i),\operatorname{dist}(g,D_i)\}
                                                        \right)\right\}.
\end{equation}
All local pieces, including the logarithmic circulation terms, cancel
exactly in this difference: the true and patch calculations have the
same retained faces and affine data near the deletion.  A physical free
wall is the same wall on both calculations.  Enlarging the fixed collars
allows the estimate for the bounded groups of simultaneous events used
next.  This is the precise reason that links have an $e^{-cR}$ factor
with no additional $\log s$ loss.

\subsection{Exact isolation and links with numerical dependencies}
\label{prep:isolation}

We now apply the finite covariant isolation construction of
\cite[Subsection 9.2]{B}.  We describe both the scalar
normalization and the support labels, because exact copying depends on
them.

Color tile centers by their tile type and the square-symmetry orbit of
their coordinates modulo $q$, in units of $R$.  In one color group nearby
centers into proximity components, with a fixed large proximity radius.
Take $q$ greater than sixteen times that radius.  A square-symmetry
orbit has at most eight residue labels.  If a group contained two
distinct centers with the same label, choose a shortest proximity path
between any such pair.  Its interior labels are distinct, so the first
repeated label occurs within eight steps.  Its displacement is therefore
less than $q$; since its endpoints have the same residue, they must be
the same center, a contradiction.  Thus each label occurs at most once
and the group has at most eight centers.
Distinct groups of the same color have disjoint enlarged neighborhoods.
Order the finite list of colors, and process groups within a color
simultaneously.  An event isolates every tile in its group, including
from the other members of the group, deleting a bond only when it is
still present.

As in the cited construction, a group's core contains its tiles and
their immediate neighbors, its open patch extends a fixed number of
tile units farther, and its guard separates core from exterior.  One
may use distances $1$, $5$, and $3$ for core, patch, and start of guard,
and distances $12$ and $10$ for an enlarged patch and guard when testing
pins.  Enlarge these fixed constants, and first the proximity radius,
whenever the fixed collars demand it.  All distances in this paragraph
are in units of $R$ observation cells.

Let $o$ and $c$ denote the current open patch before and after one such
deletion.  Its proxy is
\begin{equation}\label{prep:replacement}
 B_i^l(y)=\frac{Z_o^0(0)}{Z_c^0(0)}
              \exp\{-\tfrac12(E_o^l(y)-E_c^l(y))\}.
\end{equation}
This is \cite[Equation (9.6)]{B} with affine shape
energies.  Its constants are always centered and have no cutoff.  Delete
the physical face cutoff together with that face's microscopic bonds;
thus a cutoff is present exactly when its face is retained.

At each considered event write the original local weight as $B_i^l$
times its cut weight plus their difference.  Selecting the difference
marks that event.  After a mark skip subsequent events near that group,
with the fixed skip distance enlarged along with the collars; distance
$30$ in the preceding geometry suffices before such enlargements.
Entire groups are skipped together.  The resulting binary expansion is
an exact finite identity.  A later considered patch does not meet an
earlier marked core or guard, so its proxy is independent of the signed
choice within the earlier difference.  The fixed number of batches and
the skip rule imply that every remaining cross-tile bond lies in a
fixed tile halo of a mark.  All departures from the path with no marks
are confined to these halos.

Let $B_i^{l,(0)}$ be the proxies on the path with no marked events.  Divide
the expansion by
\begin{equation}\label{prep:base}
 \mathcal B^l(y)=\prod_iB_i^{l,(0)}(y)
                   \prod_b Z_b^0(0)e^{-E_b^l(y)/2},
\end{equation}
where $b$ runs through the isolated tiles.  This is the normalization
of \cite[Equation (9.7)]{B}, again with centered no-cutoff
constants and affine energies.  A marked halo retains its full signed
height sum and every affected proxy.  Outside these halos expand an
unused tile factor as
\begin{equation}\label{prep:tile-factor}
 \frac{e^{E_b^l(y)/2}Z_b^l(y)}{Z_b^0(0)}
          \prod_{f\in\mathcal F(b)}\chi_T(d_f)=1+F_b^l(y).
\end{equation}
In the magnetic calculation there is one exception: if the tile
containing the source remains unused, retain its whole factor rather
than expanding it around $1$.  It will be the compulsory insertion.
If the source belongs to a marked halo, that halo supplies the
compulsory factor instead.  Exactly one connected component of every
resulting pattern contains the source.

Write $E_0^l$ for the energy appearing in \eqref{prep:base}, so that
$\mathcal B^l(y)$ is a constant times $e^{-E_0^l(y)/2}$.  Combining it
with \eqref{eq:interaction} leaves the Gaussian factor
\begin{equation}\label{prep:mismatch-factor}
                    \exp\{\tfrac12(E^l(y)-E_0^l(y))\}.
\end{equation}
Telescope the energy changes along the no-mark path.  For each event,
the true change and the open-patch change have identical local parts;
their remaining difference is \eqref{prep:event-mismatch}.  Summing over
events gives a matrix $M$ on face arrays with
\begin{equation}\label{prep:link-matrix}
 E^l(y)-E_0^l(y)=\sum_{f,g}M(f,g)d_fd_g,
 \qquad |M(f,g)|\le Ce^{-cR-c\operatorname{dist}(f,g)}.
\end{equation}
Before summing, the event localization in
\eqref{prep:event-mismatch} remains available.  It makes the sum over
event positions uniformly convergent.  This argument also treats the
true free-wall clips, since both paths use the same physical wall.

\paragraph{Telescoping the numerical dependence.}
\label{prep:links}
An exponentially small coefficient computed from the whole graph is
not yet an exactly local coefficient.  We resolve it into local
calculations before assigning supports.  For each event and pair of
argument faces, take open cell boxes containing the event patch and the
two faces.  The initial radius is a sufficiently large fixed multiple
of $R$ plus such a multiple of the largest separation among the event
and argument faces.  Include a fixed $R$-size halo for every numerical
tile rule.  Increase the radius one cell at a time.  Let
$\Delta_i^{[D]}(f,g)$ denote the calculation clipped at radius $D$,
with the same grids and the same true free walls if one is met.  Write
the exact coefficient as the first clipped calculation plus successive
differences
\begin{equation}\label{prep:coefficient-telescope}
 \Delta_i(f,g)=\Delta_i^{[D_0]}(f,g)
     +\sum_{D\ge D_0}\bigl(\Delta_i^{[D+1]}(f,g)
                                  -\Delta_i^{[D]}(f,g)\bigr).
\end{equation}
On finite squares the sequence becomes exact when the clip contains the
graph.  On the plane it converges by the solution comparison in
Lemma~\ref{prep:localization}.  Every summand depends on a finite box and
has exponentially small size in its radius, after weakening constants.
The first summand has the same bound by
\eqref{prep:event-mismatch}, because its initial radius is controlled by
$R$ and the argument distances.  Subsequent differences have it by
comparison on the common interior box.  In particular the coefficients
retain both the $e^{-cR}$ gain and summable decay in the added radius.

On a torus perform the same telescoping only up to a sufficiently small
fixed fraction of the period, and then use the exact residual as a final
term.  Give that term a winding label.  If the initial radius already
reaches the threshold, no initial clipped calculation is needed; the
exact term has exponential decay in the period length because $R$ and
the argument distances then control that length.  Otherwise the last
residual has the same decay by graph comparison.  Include actual winding
paths in the support of either kind of terminal term.

Index the bilinear summands resulting from this telescoping by $\nu$,
and write $x_\nu(y)$ for one-half its coefficient times $d_fd_g$, with
the chosen symmetric matrix convention.  Define the link
$\ell_\nu(y)=e^{x_\nu(y)}-1$.  The exact version of
\cite[Equation (9.10)]{B} is now
\begin{equation}\label{prep:link-product}
          e^{(E^l(y)-E_0^l(y))/2}
                           =\prod_\nu(1+\ell_\nu(y)).
\end{equation}
At every fixed finite-volume field the coefficient series and this
product converge absolutely.  Their integrated absolute estimates are
proved in the entry section.

\paragraph{Support labels and copying.}
A link whose numerical box has radius $D'$ in cell units is covered by
a path, or a fixed number of paths, of total length $O(R+D')$.  These
paths contain its argument faces and span its numerical dependence.
Choose paths on the underlying uncut grid; deleted faces do not obstruct
a support label.  Extend them to the coordinate extrema of the
dependence box so that their spanning rectangle contains that box.
When a numerical calculation reaches a free wall, include a path to
the wall in its support.  Equivalently one may always extend the paths
to all four sides of the dependence box.  Include the symmetric image
choices to retain square covariance.  On a torus retain winding lengths
and the winding paths specified above.  The number of $m$-blocks met by
such a cover is
\begin{equation}\label{prep:link-support-size}
                             O(1+D'/R).
\end{equation}
The support is the blocks on the paths and their prescribed fixed
collars; the spanning rectangle is used to check the full region on
which its numerical coefficient depends.

These rules make copying exact.  First telescope on the same cell grids
in the plane and in a box.  Any coefficient whose dependence box and
required clearance avoid the physical wall is then the identical local
calculation in both geometries.  A coefficient that does depend on the
wall is flagged by a support reaching the wall.  On a torus, a small
unwrapped neighborhood contained in a sufficiently small fixed fraction
of the side sees precisely the plane terms, since all final residual
terms carry winding supports.  In a simply connected patch avoiding the
mark, gauge change turns the face arguments into ordinary differences;
the numerical rules are unchanged.  Near the mark the same construction
copies the pierced-plane calculation with the translation
\eqref{eq:lift-difference}.  This explains why the spanning rectangle,
including its dependence box, was required in the definition of
shift-good supports.

There are only polynomially many face, event, origin, path, and radius
labels of a given length above an anchor block.  The coefficient decay
therefore pays the exponential support weights associated with
\eqref{prep:link-support-size} when $R$ is large.  More concretely,
\[
 |\ell_\nu|\le |x_\nu|e^{|x_\nu|},
 \qquad |d_fd_g|\le\tfrac12(d_f^2+d_g^2).
\]
Assign to each link a small positive face-square cost using these
inequalities.  Summing over every radius and label gives an
$O(e^{-cR})$ total coefficient at each face.  The link fields themselves
use only their two argument faces.  No value of a branch potential along
a long covering path is estimated or inserted.  The enlarged labels
are used also for the unmarked calculation and leave the subsequent
local maps unchanged.

Finally cover every selected tile and marked event by its required halo,
and every link by the just specified paths.  Average the resulting
patterns over $e$ and the Gaussian modes below $m$ in
\eqref{eq:gaussian-observation}.  The covariance
\eqref{prep:entry-covariance} is positive and has range $O(m)$ by
\cite[Lemma 8.1]{B}; its shell continuation has the
same properties at subsequent scales.  Choose the fixed connection
radius $\rho$ large enough to include this range and the activity
collars.  Distinct connected components of a pattern then have
independent Gaussian integrations, and their numerical choices and
field factors separate.  Grouping components gives the compatible
subset algebra of \cite[Proposition 4.8]{B}.
All original exclusions are retained, including separation from the
component containing the source.  In the pierced sector exactly that
component is the compulsory insertion.  We have obtained the exact
algebraic representation and its copying rules; the remaining task is
to bound its background and insertion activities in the analytic norm.
